\documentclass[journal,lettersize]{IEEEtran}

\usepackage{aligned-overset}
\usepackage{hyperref}
\usepackage{graphicx}
\usepackage{makecell}
\usepackage{multirow}
\usepackage{tabularx}
\usepackage{multirow}
\usepackage{tabularx}
\usepackage{enumerate}
\usepackage{bbm}
\usepackage{latexsym}
\usepackage{wasysym}
\usepackage{subfigure}
\usepackage{amssymb}
\usepackage{comment}

\newcommand\noarxiv[1]{#1}

\newtheorem{thm}{\textbf{\textsc{Theorem}}}

% correct bad hyphenation here
\hyphenation{target-Origin js-rsa-sign op-tical net-works semi-conduc-tor}

\begin{document}

\title{UPPRESSO: Untraceable and Unlinkable Privacy-PREserving Single Sign-On Services}


\author{Chengqian Guo, Jingqiang Lin, Quanwei Cai, Wentian Zhu, Wei Wang, Jiwu Jing, Qiongxiao Wang, Bin Zhao, and Fengjun Li
        % <-this % stops a space
\thanks{Chengqian Guo is affiliated with Yuncheng Vocational and Technical University, China. Jingqiang Lin (corresponding author, linjq@ustc.edu.cn), Wentian Zhu, and Wei Wang are affiliated with University of Science and Technology of China (USTC), China. Quanwei Cai is with Beijing Zitiao Network Technology Co., Ltd, China. Jiwu Jing is with University of Chinese Academy of Sciences (UCAS), China. Qiongxiao Wang is with Beijing Certification Authority (BJCA) Co., Ltd, China. Bin Zhao is with Tencent America, USA. Fengjun Li is with the University of Kansas, USA.}% <-this % stops a space
\thanks{Manuscript received September 23, 2025.}
}

% The paper headers
\markboth{IEEE Transactions on Dependable and Secure Computing,~Vol.~x, No.~x, August~2021}%
{C. Guo \MakeLowercase{\textit{et al.}}: \usso: Untraceable and Unlinkable Privacy-PREserving Single Sign-On Services}

% \IEEEpubid{0000--0000/00\$00.00~\copyright~2021 IEEE}
% Remember, if you use this you must call \IEEEpubidadjcol in the second
% column for its text to clear the IEEEpubid mark.

\maketitle

% Chengqian Guo, Jingqiang Lin, Quanwei Cai, Wentian Zhu, Wei Wang, Jiwu Jing, Qiongxiao Wang, Bin Zhao, Fengjun Li

%\author{
%{\rm Chengqian Guo$^{\sharp}$, \ \ Jingqiang Lin$^{\ddag}$, \ \ Quanwei Cai$^{\P}$, \ \ Wentian Zhu$^{\ddag}$, \ \ Wei Wang$^{\ddag}$,}\\
%{\rm Jiwu Jing$^{\diamondsuit}$, \ \ Qiongxiao Wang$^{\nabla}$, \ \ Bin Zhao$^{\triangle}$, \ \ Fengjun Li$^{\S}$}\\
%$\sharp$ Yuncheng University, China\\
%$\ddag$ School of Cyber Security, University of Science and Technology of China\\
%$\P$ Beijing Zitiao Network Technology Co., Ltd, China\\
%${\diamondsuit}$ School of Computer Science \& Technology, University of Chinese Academy of Sciences\\
%${\nabla}$ Beijing Certification Authority Co., Ltd, China\ \ \ \ \ \ \ \ \ \ \ \ \ \ 
%$\triangle$ JD.com Silicon Valley R\&D Center, USA\\
%$\S$  Department of Electrical Engineering \& Computer Science, the University of Kansas, USA}



\begin{abstract}
Single sign-on (SSO) protocols such as OpenID Connect (OIDC),
allow a user to maintain only the credential for an identity provider (IdP) to log into multiple relying parties (RPs).
However, SSO introduces privacy threats, as ({\em a}) a curious IdP could trace a user's visits to RPs, and ({\em b}) colluding RPs could learn a user's online profile by linking her identities across these RPs.
This paper presents a privacy-preserving SSO scheme, called UPPRESSO, to protect a user's online profile against both (\emph{a}) an honest-but-curious IdP and (\emph{b}) malicious RPs colluding with some users.
An identity-transformation approach is proposed to generate untraceable ephemeral pseudo-identities for an RP and a user in each login, %(denoted as $PID_{RP}$ and $PID_U$, respectively)
 from which the visited RP derives a permanent account for the user, %(denoted as $Acct$),
  while these identity transformations provide unlinkability among the logins visiting different RPs.
We built the UPPRESSO prototype providing OIDC-compatible SSO services, accessed from a commercial-off-the-shelf browser without installing any extension.
%Moreover, this approach protects the identities of a user and the visited RP,
%We built a prototype of UPPRESSO on top of MITREid Connect, an open-source SSO system.
The extensive evaluations show that it fulfills the security and privacy requirements of SSO with reasonable overheads.
\end{abstract}

\maketitle

\begin{IEEEkeywords}
Single sign-on, security, privacy. %, trace, linkage
\end{IEEEkeywords}


\input{chap_Compact2/introduction.tex}
\input{chap_Compact2/background.tex}
\input{chap_Compact2/challenges.tex}
\input{chap_Compact2/threatmodel.tex}
\input{chap_Compact2/protocol.tex}
\input{chap_Compact2/analysis-no-dy.tex}
\input{chap_Compact2/implementation.tex}
\input{chap_Compact2/evaluation.tex}
\input{chap_Compact2/discussion.tex}
\input{chap_Compact2/conclusion.tex}


% conference papers do not normally have an appendix


% use section* for acknowledgement
%\section*{Acknowledgment}

% trigger a \newpage just before the given reference
% number - used to balance the columns on the last page
% adjust value as needed - may need to be readjusted if
% the document is modified later
%\IEEEtriggeratref{8}
% The "triggered" command can be changed if desired:
%\IEEEtriggercmd{\enlargethispage{-5in}}

% references section

% can use a bibliography generated by BibTeX as a .bbl file
% BibTeX documentation can be easily obtained at:
% http://www.ctan.org/tex-archive/biblio/bibtex/contrib/doc/
% The IEEEtran BibTeX style support page is at:
% http://www.michaelshell.org/tex/ieeetran/bibtex/
%\bibliographystyle{IEEEtranS}
% argument is your BibTeX string definitions and bibliography database(s)
%\bibliography{IEEEabrv,../bib/paper}
%
% <OR> manually copy in the resultant .bbl file
% set second argument of \begin to the number of references
% (used to reserve space for the reference number labels box)
%\begin{thebibliography}{1}


%\input{chap_Compact/ethics.tex}


%\end{thebibliography}
\bibliographystyle{IEEEtran}
\bibliography{ref}

%%%%%%%%%%%%%%%%%%%%%%%%%%%%%%%%%%%%%%%%%%%%%%%%%%%%%%%%%%%%%%%%%%%%%%%%%%%%%%%%%%%%%%%%%

%\noarxiv{\input{chap_Compact/appendix-rp-collision.tex}}

%\input{chap_Compact/appendix.tex}


% that's all folks
\end{document}




%几点思考：
%1. PID_RP完全由RP来决定，是不是可行？应该是不可行的，因为RP可能是恶意的，此时PID_RP并不是算出来，而是来自于另一个会话、就可以把这个Token拿到另一个会话。相当于，在我们的模型中，PID_RP = [t]ID_RP，必须是这样的一个计算步骤，否则各种证明并不成立。这个计算步骤，需要由honest程序来执行。
%相当于，恶意用户与honest RP登录（会话A），将其PID_RP拿去给另一个会话：恶意RP与honest用户的会话B；然后，把会话B的authorization code拿给会话A的恶意用户、恶意用户传给honest RP，登录结果就是有问题的：恶意用户以honest用户的名义、登录了honest RP。
%2. PID_RP不让用户知道，行不行？因为上面的攻击中，前提：需要会话A的恶意用户拿到PID_RP。
%如果用户不知道，RP单独与IdP联系、在拿到authorization code之后、RP再发送PID_RP去取Token。
%则有如下攻击：恶意用户与honest RP登录（会话A）、honest用户与恶意RP（会话B）
%终究IdP在给RP发送Token的时候，Honest RP还是应该比较一下PID_RP？否则，就是乱用了。